\section{Dependent Layered Modal Type Theory}\labeledit{sec:dt}

In this section, we combine the work by \citet{hu2024layered} and what we have built
up in the previous sections and scale all the way up to dependent types. %
We present \delamlang, \textbf{De}pendent \textbf{La}yered \textbf{M}odal Type
Theory. %
With dependent types, we can not only analyze the syntax of programs, but also that of
types. %
This ability, therefore, gives us the power to write tactics that could potentially
fill in proof obligations in a proof environment. %
In particular, this type theory addresses a number of problems that we often see in
proof assistants like Coq, Lean, and Agda. %
In Coq, tactics are written in a separate language, Ltac or Ltac2, where the
advantages of dependent types in Gallina, the core language, are lost. %
Stratifying the tactic language and the core language into two also cause
duplications: there are multiple notions for natural numbers, functions, etc.. %
On the other hand, in Lean and Agda, we use reflection to convert a Lean or an Agda
term into an AST and then use the core language to manipulate the AST. %
An instrumentation in the kernel is responsible for converting this AST back to a
valid term, if type-checked. %
This mechanism superficially provides a uniform way to tactics, but reflection
generally fails to guarantee the well-formedness of ASTs, making type malformedness
run-time errors and necessitating exception mechanisms exclusively for macro executions.

We believe that this type theory provides an example to address all aforementioned
problems. %
Starting this section, let us dive into dependent types.

\subsection{Highlights}

On a high level, we would continue to apply the layering principle in \delamlang to
guide us in the design of this type theory. %
In particular, we would want the layer for static code to be subsumed by layer for
computing programs. %
Moreover, with context variables, we are now able to formulate a recursive
elimination principle for code, which was not possible in simple types. %
However, these two requirements combining together causes some high-level technical
effects to the design of the type theory, which are worth mentioning before presenting the
type theory itself.


\subsubsection{Code Promotion}

Since we are going to introduce elimination principle for terms with dependent types,
we must also consider how equivalences are handled for code. %
For example, if we know a given piece of code has type $(\lambda x. x)~\Nat$, should
we regard this type the ``same'' as $\Nat$? %
Intuitively, the answer should be yes. %
After all, we only want to capture the syntax of the term, not its type. %
Effectively, for code of type $\cont T$, it should also be regarded as an inhabitant
of another type $\cont{T'}$, as long as $T \approx T'$ in $\Gamma$. 

In the context of dependent types, however, that causes some problems in the presence
of function applications. %
Consider a function $f : \Pi(x : S). T$ and an argument $t : S$. %
Then in general, the type of $f~t$ is $T[t/x]$. %
Now, let us construct this term as code. %
Even though $t$ is constructed as part of the code, the type of the overall code
$\cont{(T[t/x])}$ contains $t$ and therefore part of the dynamics of $t$ is in fact
promoted to the type. %
For a more concrete example, if $f : \Pi(x : \Ty 0). x$ and let the argument be
$(\lambda x. x)~\Nat$, then $\boxit{(f~((\lambda x. x)~\Nat))}$ has type
$\cont{(\lambda x. x)~\Nat}$, which we agree is just $\cont{\Nat}$. %
Clearly, the argument computes and is not purely static code as in simple types. %
We cannot avoid this phenomenon because of dependent types, so we must handle it with
care. %
This phenomenon is call a \emph{code promotion}. 

\subsubsection{Non-cumulativity}

Due to code promotions, we must permit non-trivial equivalences in types of code. %
This causes particular problems when we want an elimination principle for code with
intensional analyses. %
When we split code into cases in the elimination form, we must specify in each case
how do we construct the original code from its components. %
Therefore, it is the most convenient, when each term has a unique type, leading to a
conclusion of preferring non-cumulative universes. %
Whereas with cumulative universes, types live in higher universes for free. %
Cumulativity forms a pre-order of types which cannot be captured solely by equivalence
rules and makes the typing rule for the elimination principle extremely difficult to express, if
not completely impossible.

\subsubsection{Universe Polymorphism}

Though it is often omitted in other work, universe levels and universe polymorphism
are important ingredients in dependent type theory. %
They are typically considered as ``details'' and are not very much paid attention
to. %
However, in this work, we must be explicit about universe polymorphism. %
Consider some code for function application $t~s$. %
We in general do not know the type of $s$, let alone its universe level (though it
must be uniquely determined due to non-cumulativity discussed above). %
Therefore, the elimination principle for code must work for \emph{all} universe
levels, leading to a formalism of universe polymorphism. 

\subsubsection{Tarski-style Universes}

Another ingredient to consider when approaching an elimination principle for code with
dependent types and intensional analyses is the separation between types and terms. %
Consider Russell-style universes where types and terms are not distinguished
naturally. %
It would probably suffice to say that $\cont T$ can represent code for both types and
terms. %
In particular, code for some types just has type $\cont{\Ty 0}$, for example. %
Unfortunately, this thought is too naive. %
When we consider $\cont T$ as the type of code, we are considering this type with two
indices, the (local) context $\Gamma$ and the type of the code body $T$. %
But what is the type of $T$? %
Well, it is $\cont{\Ty l}$ for some $l$, which is just a special case of $\cont T$! %
A type clearly should not be indexed by a special case of itself. %
The intertwine between types and terms in Russell-style universes seems to even
prevent a proper statement of indices of types for code. %
However, Tarski-style universes, where types and terms are clearly distinguished, introduce
mutually inductive relations between types and terms, and safely bail us out of this
problem, as we will see very soon.

\subsection{Syntax}\labeledit{sec:dt:syntax}

Let us start with the syntax of \delamlang. %
Since we employ Tarski-style universes, the syntax of terms and types are separate. %
Due to non-cumulativity, certain constructs must remember universe levels. %
Due to the elimination principle of code, some constructs must include additional
sub-structures so that the elimination eventually checks out. %
Let us begin with the subset that is basically just Martin-L\"of type theory (MLTT).
\begin{alignat*}{2}
  x, y & && \tag{Local variables} \\
  \ell & && \tag{Universe variables} \\
  l & := &&\ \ell \sep 0 \sep 1 + l \sep l \sqcup l' \sep \omega \tag{Universe levels} \\
  M, S, T & := && \ \Nat \sep \PI{l}{l'}x S T \sep \Ty l \sep \UPI \ell l T \sep \Elt l t \tag{Types, $\Typ$} \\
  s, t &:=&&\ x \tag{Terms, $\Exp$} \\
  & && \sep \Nat \sep \PI{l}{l'}x s t \sep \Ty l \tag{encoding of types}  \\
  & && \sep \ze \sep \su t \sep \ELIMN l{x.M}{s}{x, y.s'}{t} \tag{natural numbers $N$} \\
       & &&\sep \LAM l{l'} x S t \sep \APP t l{l'} x S T s \tag{dependent functions}  \\
  & && \sep \ULAM l \ell t \sep \UAPP t l \tag{universe polymorphic functions} \\
  \Gamma, \Delta &:= &&\ \cdot \sep g \sep \Gamma, x : T \at l
                        \tag{Local contexts, \Ctx} \\
  L &:= && \cdot \sep L, \ell \tag{Universe contexts}
\end{alignat*}
Due to three kinds of contexts and the scale of the system, we omit the discussion of
weakenings, which we diligently kept track of in the previous sections. %
We take various weakenings for granted. %
Nevertheless, they will appear in the semantics. %
Weakenings in general are pretty obvious, as we permit arbitrary lookups for all
variables. %
Since we now must deal with universe levels, we use $\ell$ to range over variables for
universe levels. %
The syntax for universe levels follows Agda's conventions. %
Universe levels form an idempotent commutative monoid, the laws of which we will show
in the next subsection. %
Here we use $\sqcup$ to denote taking the max of two universe levels. %
% Our syntax above has not introduced the syntax to introduce a universe level
% variables.
The ability to take maximum between two levels induces a partial order:
\begin{align*}
  l \le l' := l' \approx l \sqcup l'
\end{align*}
where we use $\approx$ to express the equivalence between universe levels. %
Thus, with $\le$, universe levels form a bounded inf-lattice. %
A strict order is given by requiring the pre-order to hold for the successor of $l$:
\begin{align*}
  l < l' := l' \approx (1 + l) \sqcup l'
\end{align*}
This strict order, as to be shown later, is well-founded, based on which we will give
semantics to universes. %
Due to universe polymorphism, we must also include an $\omega$ level, which will be
used to represent the universe level of a universe-polymorphic function. %
The $\omega$ level must not appear in any program, does not participate in the bounded
inf-lattice specified above, is only used in type-checking, and therefore can be
ignored most of the time. %
The formalization of universe polymorphism here follows~\citet{bezem_type_2023}
tightly. %
We use $\UPI \ell l T$ to denote the type of a universe-polymorphic function. %
It introduce a non-empty list of universe variables at once, and lives at $\Ty
\omega$. %
The type $T$ lives at $\Ty l$, where $l$ may refer to all variables from $\vect
\ell$. %
Since $l$ cannot be $\omega$, we must have all universe variables introduced in one
go. %
The introduction form is $\ULAM l \ell t$, which also introduces a non-empty list of
universe variables first and then the function body as expected. %
The elimination form $\UAPP t l$ symmetrically eliminates a universe-polymorphic
function with the same number of universe level expressions. %

The rest of the expressions are pretty much standard from MLTT. %
We have natural numbers (\Nat), their introduction forms and a recursion principle. %
We always use $M$ to exclusively represent the motives of a recursion principle. %
For regular dependent functions $\PI{l}{l'}x S T$, we must specify the universe levels
of $S$ and $T$, following~\citet{pujet_impredicative_2023}. %
We might omit the universe levels if they are not important in the discussion. %
The function abstraction $\LAM l{l'} x S t$ is standard; we might omit $l$ and $S$ if they
are not important or can be inferred from the context. %
The function application $\APP t l{l'} x S T s$ is arguably more complex. %
We explicitly specify the type of the function to prepare for a better formulation of
the elimination principle for code later in the section. %
By requiring explicit type annotations in elimination forms, types that are usually
hidden in the core syntax become 
sub-structures in the elimination form for code. %
This verbosity has no negative impact for programmers: after all, we are discussing a core theory, and
we can let a type-inferring front-end to fill in these types for the users if they
choose so. %
Following conventions, we may simply write $s~t$ if the types are not important. %

Since we are employing Tarski-style universes, as we have specified in the syntax,
types and terms are separated. %
As terms, we use \emph{encodings} of types, i.e. the overloaded $\Nat$,
$\PI l{l'} x s t$ and $\Ty l$, which are members of some universes. %
They are \emph{decoded} into actual types through $\Elt l t$, converting the encodings
to actual corresponding types. %
This part is basically identical to~\citet{palmgren_universes_1998}.  %
For simplicity, we have omitted the type lifter, which is responsible for raising the
universe levels explicitly, similar to \texttt{Lift} in Agda's standard library. %
According~\citet{palmgren_universes_1998}, the type lifter bears additional
equivalences and therefore we omit them here for conciseness, as lifting of the levels
is an orthogonal issue here. %

\subsection{Universe Levels}\labeledit{sec:dt:ulevel}

In the syntax, we deliberately group all the universe variables into a separate
context. %
This is beneficial as both local and global contexts (to be discussed later) will need to refer
to universe variables. %
It is also helpful for the future work of extending \delamlang to more
layers, by simply inserting more contexts after $L$. %
In this section, we state the well-formedness and equivalence judgments for universe
levels. %
Note that all syntactically valid universe contexts are already well-formed as they
only contain universe variables.
\begin{mathpar}
  \inferrule
  {\ell \in L}
  {\typing[L]{\ell}{\Level}}

  \inferrule
  { }
  {\typing[L]{0}{\Level}}

  \inferrule
  {\typing[L]{l}{\Level}}
  {\typing[L]{1 + l}{\Level}}
  
  \inferrule
  {\typing[L]{l}{\Level} \\ \typing[L]{l'}{\Level}}
  {\typing[L]{l \sqcup l'}{\Level}}
\end{mathpar}
Notice that $\omega$ is not well-formed. %
Indeed, the judgment $\typing[L]l\Level$ only captures the well-formed universe levels
that can be written by a programmer. %
The level $\omega$, on the other hand, only appears during type-checking to denote the
type of universe-polymorphic functions. %

As discussed above, universe levels themselves form an idempotent commutative
monoid. %
Hence they have the following equivalence rules:
\begin{mathpar}
  \inferrule
  {\typing[L]{l}{\Level}}
  {\tyequiv[L]{l}{l}\Level}
  
  \inferrule
  {\tyequiv[L]{l}{l'}{\Level}}
  {\tyequiv[L]{l'}{l}\Level}

  \inferrule
  {\tyequiv[L]{l}{l'}\Level \\ \tyequiv[L]{l'}{l''}\Level}
  {\tyequiv[L]{l}{l''}\Level}
\end{mathpar}
First we specify the basic PER rules. %
Then we have congruence rules:
\begin{mathpar}
  \inferrule
  {\tyequiv[L]{l}{l'}{\Level}}
  {\tyequiv[L]{1 + l}{1 + {l'}}\Level}

  \inferrule
  {\tyequiv[L]{l}{l'}\Level \\ \tyequiv[L]{l''}{l'''}\Level}
  {\tyequiv[L]{l \sqcup l''}{l' \sqcup l'''}\Level}  
\end{mathpar}
Finally we have the algebraic rules:
\begin{mathpar}
  \inferrule
  {\typing[L]{l}{\Level}}
  {\tyequiv[L]{l \sqcup 0}{l}\Level}
  
  \inferrule
  {\typing[L]{l}{\Level} \\ \typing[L]{l'}{\Level} \\ \typing[L]{l''}{\Level}}
  {\tyequiv[L]{(l \sqcup l') \sqcup l''}{l \sqcup (l' \sqcup l'')}\Level}

  \inferrule
  {\typing[L]{l}{\Level} \\ \typing[L]{l'}{\Level}}
  {\tyequiv[L]{l \sqcup l'}{l' \sqcup l}\Level}

  \inferrule
  {\typing[L]{l}{\Level}}
  {\tyequiv[L]{l \sqcup l}{l}\Level}

  \inferrule
  {\typing[L]{l}{\Level} \\ \typing[L]{l'}{\Level}}
  {\tyequiv[L]{1 + (l \sqcup l')}{1 + l \sqcup 1 + {l'}}\Level}

  \inferrule
  {\ell \in L}
  {\tyequiv[L]{\ell \sqcup (1 + \ell)}{1 + \ell}{\Level}}
\end{mathpar}
The second last rule is distributivity of $\tsucc$ over $\sqcup$. %
The last rule is absorption of $\tsucc$ over $\sqcup$. %
The equivalence judgment confirms the well-formedness of both components:
\begin{lemma}[Presupposition]
  If $\tyequiv[L]{l}{l'}\Level$, then $\typing[L]l\Level$ and $\typing[L]{l'}\Level$. 
\end{lemma}
\begin{proof}
  Induction.
\end{proof}
We can prove a more general absorption rule by doing a few inductions.
\begin{lemma}[Absorption]
  If $\typing[L]l\Level$, then $\tyequiv[L]{l \sqcup (1 + l)}{1 + l}\Level$. 
\end{lemma}
\begin{proof}
  Induction. Only the following case is interesting:
  \begin{mathpar}
    \inferrule
    {\typing[L]{l}{\Level} \\ \typing[L]{l'}{\Level}}
    {\typing[L]{l \sqcup l'}{\Level}}
  \end{mathpar}
  We reason as follows:
  \begin{align*}
    l \sqcup l' \sqcup (1 +{l \sqcup l'})
    & \approx (l \sqcup (1 + l)) \sqcup (l' \sqcup (1 +{l'})) \\
    & \approx (1 + l) \sqcup (1 +{l'}) \byIH \\
    & \approx 1 +{(l \sqcup l')}
  \end{align*}
  Hence the proof is complete. 
\end{proof}
Then we generalize further:
\begin{lemma}[Absorption]
  If $\typing[L]l\Level$, then for any $n$, $\tyequiv[L]{l \sqcup (n + {l})}{n + {l}}\Level$.   
\end{lemma}
\begin{proof}
  We proceed by induction on $n$. %
  The cases for $n=0$ and $n=1$ are simple.  We consider the step case, where
  $n = 1 + n'$ and we know
  $\tyequiv[L]{l \sqcup (n' + {l})}{n' + {l}}\Level$. %
  We reason as follows:
  \begin{align*}
    l \sqcup (1 + (n' + l))
    & \approx l \sqcup (1 +{(l \sqcup (n' + {l}))})
      \byIH \\
    & \approx l \sqcup (1 + l) \sqcup (1 + n' + {l}) \\
    & \approx (1 + l) \sqcup (1 + n' + {l})
    \tag{by absorption} \\
    & \approx (1 + l) \sqcup (n' + (1 + l)) \\
    & \approx n' + {(1 + l)} \byIH \\
    & \approx n + l
  \end{align*}
  The proof is complete. 
\end{proof}

As readers might have noticed, the theory for universe levels are self-contained and
their equivalence is decidable, as per implemented by Agda's type-checker. %
For this reason, in the remainder of the discussion, we undermine the importance of
well-formedness and equivalence of universe levels, unless it is essential in the
surrounding context. 

Next, we define substitutions for universe levels:
\begin{align*}
  \phi := \cdot \sep \phi,l/\ell \tag{Substitutions for universe levels}
\end{align*}
\begin{mathpar}
  \inferrule
  { }
  {\typing[L]{\cdot}{\cdot}}

  \inferrule
  {\typing[L]\phi{L'} \\ \typing[L]l\Level}
  {\typing[L]{\phi, l/\ell}{L',\ell}}
\end{mathpar}
Applying substitutions is intuitive:
\begin{align*}
  \ell [\phi] &:= \phi(\ell) \\
  0 [\phi] &:= 0 \\
  1 + l [\phi] &:= 1 + (l[\phi]) \\
  l \sqcup l' [\phi] &:= (l [\phi]) \sqcup (l'[\phi]) \\
  \omega [\phi] &:= \omega
\end{align*}
We then have the following lemmas:
\begin{lemma}
  If $\typing[L]l\Level$ and $\typing[L']\phi L$, then $\typing[L']{l[\phi]}\Level$.
\end{lemma}
\begin{proof}
  Induction.
\end{proof}
\begin{lemma}
  If $\tyequiv[L]l{l'}\Level$ and $\typing[L']\phi L$, then $\tyequiv[L']{l[\phi]}{l'[\phi]}\Level$.
\end{lemma}
\begin{proof}
  Induction.
\end{proof}

\newcommand{\adjustf}[1]{\ensuremath{\textsf{adjust}(#1)}}

The well-foundedness of the strict order $<$ is intuitive. %
The only bottom element is $0$. %
We simply keep removing $\tsucc$ from all components of $\sqcup$, and we must
eventually stop. %
Thus, the simplest way to argue the well-foundedness of $<$ is to define a measure
based on the number of $\tsucc$'s that can be removed from an $l$. %
This number is defined over all $\typing[L]l\Level$ recursively as follows:
\begin{align*}
  \countf \ell &:= \{ \ell \mapsto 0, 0 \mapsto 0 \} \\
  \countf 0 &:= \{ 0 \mapsto 0 \} \\
  \countf{1 + l} &:= \mergef{\{ \ell \mapsto 1 + n \sep \ell \mapsto n \in \countf{l}
                   \}}{\{0 \mapsto 1 + \countf{l}(0)\}} \\
  \countf{l \sqcup l'} &:= \mergef{\countf{l}}{\countf{l'}}
\end{align*}
Here \textsf{count} returns a map that counts the number of $\tsucc$'s over all
universe variables and $0$. %
The function \textsf{merge} merges two maps and takes the maximum in a conflict. %
The following definition makes sure equivalent universe levels to have the same
representation as maps:
\begin{align*}
  \adjustf m :=
  \begin{cases}
    \{ \ell \mapsto n \sep m(\ell) \mapsto n \} & \text{if $m(0) \le \max_{\ell
                                                  \mapsto n \in m} n$} \\
    m & \text{otherwise}
  \end{cases}
\end{align*}
In the first branch, we check if there is a variable which has a higher universe level
than the constant. %
If so, we drop the constant completely. %
For example, in $2 \sqcup (2 + \ell)$, $2$ is
redundant, as we know $2 + \ell$ is at least as large as $2$. %
On the other hand, in $2 \sqcup (1 + \ell)$, it is possible for $1 + \ell$
to be smaller than $2$ when we take $\ell$ as $0$. %
Therefore, in this case, we must keep $2$.  Thus, the finiteness of
decreasing steps of the universe levels can be seen as taking some finite steps by
removing all $1 +$ from maps returned by $\adjustf{\countf l}$. %
Then we can just take away variables until we can no longer descend. %
In fact, $\adjustf{\countf l}$ should be considered a normalization algorithm for
universe levels. %
We can simply compare equality between maps computed as such decide whether two
universe levels are equivalent. %
The correctness is as follows:
\begin{lemma}
  If $\tyequiv[L]l{l'}\Level$, then $\adjustf{\countf l} = \adjustf{\countf{l'}}$.
\end{lemma}
\begin{proof}
  Induction. Take advantage of the idempotent commutative monoidal nature of maximum. 
\end{proof}
This lemma ensures that the procedure respects equivalence between universe levels. %
The other direction is seen by providing a ``decoding'' function, which converts a map
to a universe level. %
% We then simply need to agree on a serialization of the map which is a mere detail
% here. %
We give one possible function that converts a map to a $\Level$.
\begin{align*}
  \flatten L m := \begin{cases}
    \bigsqcup_{\ell \mapsto n \in m} n + \ell & \text{if $0$ is not in $m$} \\
    m(0) \sqcup (\bigsqcup_{\ell \mapsto n \in m} n + \ell)
    & \text{otherwise}
  \end{cases}
\end{align*}
where the order of $\ell$'s respects their order in $L$ and all $\sqcup$ are right
associative. %
These requirements give a syntactically unique flattening of a map. %
Then we prove
\begin{lemma}
  If $\typing[L]l\Level$, then $\tyequiv l{\flatten{L}{\adjustf{\countf l}}}\Level$. 
\end{lemma}
\begin{proof}
  We proceed by induction. %
  It is rather immediate. For the $\tsucc$ case, we use its distributivity to
  propagate it inwards. For the $\sqcup$ case, we use commutativity to rearrange levels
  within and absorption to eliminate small levels whenever necessary.
\end{proof}
% We can in fact just compare $\Level$'s by considering the results of \textsf{count}. %
% This is pretty simple so we often consider the equalities of our types and terms
% modulo the equivalence of $\Level$s. 

\subsection{Typing and Equivalence Judgments}\labeledit{sec:dt:mltt-rules}

In this section, we introduce the typing and equivalence judgments, only for the MLTT
part. %
We will consider the modal part next altogether. %
The typing and equivalence judgments are defined mutually as usual. %
All the related judgments are:
\begin{itemize}
\item $\judge[L]\Psi$ denotes the well-formedness of the global context $\Psi$ under
  $L$. 
\item $\lpjudge i \Gamma$ denotes the well-formedness of $\Gamma$ given the universe
  context $L$ and the global context $\Psi$ at layer $i$. %
  In this section, we are not very concerned about layers yet as most parts about
  meta-programming and intensional analysis come in a later section
  (\Cref{sec:dt:modal}). 
\item $\lpequiv i \Gamma\Delta$ denotes the equivalence between two local contexts
  $\Gamma$ and $\Delta$. 
\item $\lttypwf i T l$ denotes the well-formedness of the type $T$ living in the
  universe level $l$ at layer $i$ in the given contexts. 
\item $\lttyping i t T l$ denotes the well-typedness of $t$ of type $T$, which is in
  the universe level $l$ at layer $i$.
  In the special occasion of $T$ being some $\Se$, we might write $\lttypingd i T
  l{1 + l}$ to simultaneously denote two judgments at the same time to save space.
\item $\lttypeq i{T}{T'}l$ denotes the equivalence between types $T$ and $T'$ living in the
  universe level $l$ at layer $i$ in the given contexts. 
\item $\lttyequiv i t {t'} T l$ denotes the equivalence between $t$ and $t'$ of type
  $T$, which is in the universe level $l$ at layer $i$. %
  The shorthand $\lttyequivd i T {T'} l {1 + l}$ has a meaning similar to above.
\item $\ltsubst i \delta \Delta$ denotes the well-formedness of a local substitution
  $\delta$ which substitute all local variables in $\Delta$ into terms referring to
  $\Gamma$. %
  We will introduce this judgment and the next when we discuss the modal part. %
\item $\ltsubeq i\delta{\delta'}\Delta$ denotes the equivalence between local
  substitutions $\delta$ and $\delta'$. 
% \item $\lttypingv i t T l$ is a sub-judgment of $\lttyping i t T l$ which makes sure
%   $t$ is a local variable or a global variable that can only be substituted to a local
%   variable.
% \item $\ltsubstv i \delta\Delta$ is the extension of $\lttypingv i t T l$ to local substitution.
\end{itemize}

The judgments for MLTT are rather routine. %
Many are just generalization of the judgments in~\Cref{sec:cv}. 
Let us first consider the well-formedness and equivalence of local contexts:
\begin{mathpar}
  \inferrule
  {\judge[L] \Psi}
  {\lpjudge i \cdot \\ \lpequiv i \cdot\cdot}

  \inferrule
  {\judge[L] \Psi \\ g : \Ctx \in \Psi}
  {\lpjudge i{g} \\ \lpequiv i{g}{g}}

  \inferrule
  {\lpjudge i \Gamma \\\\ \lttypwf i T l \\ \typing[L]l\Level}
  {\lpjudge i{\Gamma, x : T \at l}}

  \inferrule
  {\lpequiv i \Gamma\Delta \\ \lttypeq i T{T'} l \\ \lttypeq[\Psi][\Delta] i T{T'} l
    \\ \lttypwf i T l \\ \lttypwf[\Psi][\Delta] i{T'}{l'} \\ \tyequiv[L]l{l'}\Level}
  {\lpequiv i{\Gamma, x : T \at l}{\Delta, x : T' \at{l'}}}
\end{mathpar}
The well-formedness of types are also immediate, following~\citet{pujet_impredicative_2023}. %
When we encounter $\tEl$, we resort that to the typing judgment of terms. %
Overlapping rules for well-typedness of encoding as terms are also listed:
\begin{mathpar}
  \inferrule
  {\lpjudge{\typeof i} \Gamma}
  {\lttypingd i \Nat 0 {1}}

  \inferrule
  {\typing[L]l\Level  \\ \typing[L]{l'}\Level \\\\
    \lttypwf i S l \\ \lttypwf[\Psi][\Gamma, x : S \at l] i T{l'}}
  {\lttypwf i{\PI l{l'} x S T}{l \sqcup l'}}

  \inferrule
  {\typing[L]l\Level  \\ \typing[L]{l'}\Level \\\\
    \lttyping i s{\Ty l}{1 + l} \\ \lttyping[\Psi][\Gamma, x : \Elt l s \at l] i t{\Ty{l'}}{1 + l'}}
  {\lttyping i{\PI l{l'} x s t}{\Ty{l \sqcup l'}}{1 + {(l \sqcup l')}}}

  \inferrule
  {\typing[L] l \Level \\\\ \lttyping i t{\Ty l}{1 + l}}
  {\lttypwf i{\Elt l t}{l}}

  \inferrule
  {\lpjudge{\typeof i} \Gamma \\ \typing[L]l\Level}
  {\lttypingd i{\Ty l}{1 + l}{2 + l}}

  \inferrule
  {\lpjudge \metalevel \Gamma \\ \lttypwf[\Psi][\Gamma][L,\vect\ell] 1 T l \\\\ |\vect \ell| > 0 \\ \typing[L,\vect\ell]l\Level}
  {\lttypwf \metalevel {\UPI \ell l T}{\omega}}

  \inferrule
  {\lttypwf i T{l'} \\\\ \tyequiv[L]l{l'}\Level}
  {\lttypwf i{T}{l}}
\end{mathpar}
Note that all type constructors with explicitly specified universe levels must not
refer to $\omega$. %
Indeed, $\omega$ level only appears when we validate a universe-polymorphic function
and nowhere else. %
Nor can we pass around a universe-polymorphic function. %
Moreover, universe-polymorphic functions are only available at the highest layer, which is the
layer with capability to do meta-programming and recursive intensional analysis. %
This is because that universe variables must also be visible by the bindings in global
contexts. %
% In general, universe-polymorphic functions must only be available at the highest
% layer, if \delamlang is extended to more than two layers. % 
In the well-formedness rule for $\Nat$, we use a function $\typeof i$ which alters the
layer, in which local contexts live. %
This treatment is necessary to accommodate code promotions and permit computation in the
local contexts and on the type level. %
We give the actual definition of $\typeof i$ in \Cref{sec:dt:modal}.

The equivalence between types is composed of three parts. %
The first part is the PER rules. %
\begin{mathpar}
  \inferrule
  {\lttypeq i{T}{T'}l}
  {\lttypeq i{T'}{T}l}

  \inferrule
  {\lttypeq i{T}{T'}l \\ \lttypeq i{T'}{T''}l}
  {\lttypeq i{T}{T''}l}
\end{mathpar}
Then we have the congruence rules, which simply straightforwardly propagate
equivalence downwards:
\begin{mathpar}
  \inferrule
  {\lpjudge{\typeof i} \Gamma}
  {\lttyequivd i \Nat \Nat 0 {1}}

  \inferrule
  {\tyequiv[L]{l_1}{l_3}\Level \\ \tyequiv[L]{l_2}{l_4}\Level \\\\ \lttypeq i S{S'}{l_1} \\ \lttypeq[\Psi][\Gamma, x : S \at{l_1}] iT{T'}{l_2}}
  {\lttypeq i{\PI{l_1}{l_2} x S T}{\PI{l_3}{l_4} x{S'}{T'}}{l_1 \sqcup l_2}}

  \inferrule
  {\tyequiv[L]{l_1}{l_3}\Level \\ \tyequiv[L]{l_2}{l_4}\Level \\\\
    \lttyequiv i s{s'}{\Ty{l_1}}{1 + l_1} \\
    \lttyequiv[\Psi][\Gamma, x : \Elt{l_1} s \at{l_1}] i t{t'}{\Ty{l_2}}{1 + l_2}}
  {\lttyequiv i{\PI{l_1}{l_2} x s t}{\PI{l_3}{l_4} x{s'}{t'}}{\Ty{l_1 \sqcup l_2}}{1 + {(l_1 \sqcup l_2)}}}

  \inferrule
  {\tyequiv[L]{l}{l'}\Level \\ \lttyequiv i t{t'}{\Ty l}{1 + l}}
  {\lttypeq i{\Elt l t}{\Elt{l'}{t'}}{l}}

  \inferrule
  {\lpjudge{\typeof i}\Gamma \\ \tyequiv[L]l{l'}\Level}
  {\lttyequivd i{\Ty l}{\Ty{l'}}{1 + l}{2 + l}}

  \inferrule
  {\lttypwf[\Psi][\Gamma][L,\vect\ell] \metalevel T l \\
    \lttypeq[\Psi][\Gamma][L,\vect\ell] \metalevel T{T'} l \\\\ |\vect \ell| > 0 \\ \tyequiv[L,\vect\ell]l{l'}\Level}
  {\lttypeq \metalevel {\UPI \ell l T}{\UPI \ell{l'}{T'}}{\omega}}

  \inferrule
  {\lttypeq i T{T'}{l'} \\ \tyequiv[L]l{l'}\Level}
  {\lttypeq i{T}{T'}{l}}
\end{mathpar}

Finally, we have a number of computation rules that decode terms into types:
\begin{mathpar}
  \inferrule
  {\lpjudge{\typeof i} \Gamma}
  {\lttypeq i \Nat{\Elt 0 \Nat} 0}

  \inferrule
  {\typing[L]l\Level \\ \lpjudge{\typeof i} \Gamma}
  {\lttypeq i{\Ty l}{\Elt{1 + l}{\Ty l}}{1 + l}}
  
  \inferrule
  {\typing[L]l\Level  \\ \typing[L]{l'}\Level \\
    \lttyping i s{\Ty l}{1 + l} \\ \lttyping[\Psi][\Gamma, x : \Elt l s \at l] i t{\Ty{l'}}{1 + l'}}
  {\lttypeq i{\PI l{l'} x{\Elt l s}{\Elt{l'} t}}{\Elt{l \sqcup l'}{\PI l{l'} x s t}}{l \sqcup l'}}
\end{mathpar}
We do not have an encoding for universe-polymorphic functions, so there is not a
decoding rule for them.

Then we move on to defining the typing rules for terms. %
They are pretty much straightforward:
\begin{mathpar}
  \inferrule
  {\lpjudge{\typeof i}\Gamma \\ x : T \at l \in \Gamma}
  {\lttyping i x T l}

  \inferrule
  {\lpjudge{\typeof i}\Gamma}
  {\lttyping i \ze \Nat 0}

  \inferrule
  {\lttyping i t \Nat 0}
  {\lttyping i {\su t} \Nat 0}

  \inferrule
  {\typing[L]l\Level \\ \lttypwf[\Psi][\Gamma, x : \Nat \at 0] i M l \\
    \lttyping i s {M[\ze/x]}l \\
    \lttyping[\Psi][\Gamma, x : \Nat \at 0, y : M \at l] i {s'}{M[\su x/x]}l \\
  \lttyping i t \Nat 0}
  {\lttyping i{\ELIMN l{x.M}s{x,y. s'}t}{M[t/x]}{l}}

  \inferrule
  {\typing[L]l\Level  \\ \typing[L]{l'}\Level \\ \lttypwf i S l \\ \lttyping[\Psi][\Gamma, x : S \at l] i t{T}{l'}}
  {\lttyping i{\LAM l{l'} x S t}{\PI l{l'} x S T}{l \sqcup l'}}

  \inferrule
  {\typing[L]l\Level  \\ \typing[L]{l'}\Level \\ \lttypwf i S l \\ \lttypwf[\Psi][\Gamma, x : S \at l] i{T}{l'} \\
    \lttyping i t{\PI l{l'} x S T}{l \sqcup l'} \\ \lttyping i s S l}
  {\lttyping i{\APP t l{l'} x S T s}{T[s/x]}{l'}}

  \inferrule
  {\lttyping[\Psi][\Gamma][L, \vect \ell]\metalevel{t}{T}{l} \\\\ |\vect \ell| > 0 \\ \typing[L,\vect\ell]l\Level}
  {\lttyping \metalevel {\ULAM l \ell t}{\UPI \ell l T}{\omega}}

  \inferrule
  {\lttypwf[\Psi][\Gamma][L,\vect\ell] \metalevel T l \\
    \lttyping \metalevel{t}{\UPI \ell l T}{\omega} \\\\ |\vect\ell| = |\vect l| > 0 \\ \forall l' \in \vect l ~.~ \typing[L]{l'}\Level}
  {\lttyping \metalevel{\UAPP t l}{T[\vect l/\vect \ell]}{l[\vect l/\vect \ell]}}

  \inferrule
  {\lttyping i t{T'}l \\ \lttypeq{\typeof i}T{T'}l}
  {\lttyping i t T l}

  \inferrule
  {\lttyping i tT{l'} \\ \tyequiv[L]l{l'}\Level}
  {\lttyping i t T l}
\end{mathpar}

The equivalence rules for terms are also composed of three parts. %
The PER rules are immediate:
\begin{mathpar}
  \inferrule
  {\lttyequiv i t{t'}T l}
  {\lttyequiv i{t'}tT l}

  \inferrule
  {\lttyequiv i t{t'}T l \\ \lttyequiv i{t'}{t''}T l}
  {\lttyequiv i t{t''}T l}
\end{mathpar}
The congruence rules are naturally induced by the typing rules:
\begin{mathpar}
  \inferrule
  {\lpjudge{\typeof i}\Gamma \\ x : T \at l \in \Gamma}
  {\lttyequiv i x x T l}

  \inferrule
  {\lpjudge{\typeof i}\Gamma}
  {\lttyequiv i \ze \ze \Nat 0}

  \inferrule
  {\lttyequiv i t{t'} \Nat 0}
  {\lttyequiv i {\su t}{\su{t'}} \Nat 0}

  \inferrule
  {\tyequiv[L]l{l'}\Level \\ \lttypeq[\Psi][\Gamma, x : \Nat \at 0] i M{M'} l \\
    \lttyequiv i {s_1}{s_3} {M[\ze/x]}l \\
    \lttyequiv[\Psi][\Gamma, x : \Nat \at 0, y : M \at l] i{s_2}{s_4}{M[\su x/x]}l \\
  \lttyequiv i t{t'} \Nat 0}
  {\lttyequiv i{\ELIMN l{x.M}{s_1}{x,y. s_2}t}{\ELIMN{l'}{x.M'}{s_3}{x,y. s_4}{t'}}{M[t/x]}{l}}

  \inferrule
  {\tyequiv[L]{l_1}{l_3}\Level  \\ \tyequiv[L]{l_2}{l_4}\Level \\ \lttypwf i S{l_1} \\ \lttypeq i
    S{S'}{l_1} \\
    \lttyequiv[\Psi][\Gamma, x : S \at {l_1}] i t{t'}{T}{l_2}}
  {\lttyequiv i{\LAM {l_1}{l_2} x S t}{\LAM {l_3}{l_4} x{S'}{t'}}{\PI l{l'} x S T}{l_1 \sqcup l_2}}

  \inferrule
  {\tyequiv[L]{l_1}{l_3}\Level  \\ \tyequiv[L]{l_2}{l_4}\Level \\ \lttypwf i S{l_1} \\ \lttypeq i S{S'}{l_1}
    \\
    \lttypeq[\Psi][\Gamma, x : S \at{l_1}] i{T}{T'}{l_2} \\
    \lttyequiv i t{t'}{\PI {l_1}{l_2} x S T}{l_1 \sqcup l_2} \\ \lttyequiv i s{s'} S{l_1}}
  {\lttyequiv i{\APP t {l_1}{l_2} x S T s}{\APP{t'} {l_3}{l_4} x {S'}{T'}{s'}}{T[s/x]}{l_2}}

  \inferrule
  {\lttyequiv[\Psi][\Gamma][L, \vect \ell]\metalevel{t}{t'}{T}{l} \\ |\vect \ell| > 0 \\ \tyequiv[L,\vect\ell]l{l'}\Level}
  {\lttyequiv \metalevel {\ULAM l \ell t}{\ULAM {l'} \ell{t'}}{\UPI \ell l T}{\omega}}

  \inferrule
  {\lttypwf[\Psi][\Gamma][L,\vect\ell] \metalevel T l \\
    \lttyequiv \metalevel{t}{t'}{\UPI \ell l T}{\omega} \\ |\vect\ell| = |\vect l| = |\vect
    l'| > 0 \\
    \forall 0 \le n < |\vect l| ~.~ \tyequiv[L]{\vect l(n)}{\vect l'(n)}\Level}
  {\lttyequiv \metalevel{\UAPP t l}{t'~\$~\vect l'}{T[\vect l/\vect \ell]}{l[\vect l/\vect \ell]}}

  \inferrule
  {\lttyequiv i t{t'}{T'}l \\ \lttypeq{\typeof i}T{T'}l}
  {\lttyequiv i t {t'} T l}

  \inferrule
  {\lttyequiv i t{t'}T{l'} \\ \tyequiv[L]l{l'}\Level}
  {\lttyequiv i t{t'}T l}
\end{mathpar}

\subsection{Meta-programming Modalities}\labeledit{sec:dt:modal}

In this part, we introduce the modalities for meta-programming and intensional
analysis. %
We use the $\square$ modality to represent the type of
code and we use layers to control the computational behaviors of the type theory. %
However, there are two points that we need to pay attention to:
\begin{enumerate}
\item Previously, we have discussed code promotions. %
  Code promotions imply that we must permit computation of code on the type level. %
  This further implies that we must introduce an intermediate layer between layer for code and
  that for meta-programs, which restricts the language to still MLTT but permits computation. 
\item Due to Tarski-style universes, we must introduce two kinds of contextual types,
  one for types and one for terms. %
  As seen in the syntax of types and terms, they are mutually defined, so the
  recursive principles for code of types and terms must also be mutual.
\end{enumerate}

\begin{table}[]
\begin{tabular}{|l|l|l|l|l|}
\hline
Layer            & $\varlevel$ & $\codelevel$  & $\proglevel$  & $\metalevel$                          \\ \hline
Language         & Variables only &MLTT & MLTT & MLTT extended with meta-programming \\ \hline
Computation      & No & No   & Yes  & Yes                          \\ \hline
Meta-programming & No & No   & No  & Yes                          \\ \hline
Layer of types   & $\proglevel$ & $\proglevel$  & $\proglevel$  & $\metalevel$                          \\ \hline
\end{tabular}
\caption{Features at each layer}\labeledit{tab:dt:feat}
\end{table}

Having set up the basic theme, let us begin with the syntax of the extension to MLTT:
\begin{alignat*}{2}
  i &:= &&\ \varlevel \sep \codelevel \sep \proglevel \sep \metalevel \tag{Layers} \\
  U & && \tag{Global variables as types} \\
  u & && \tag{Global variables as terms} \\
  % u_\# & && \tag{Global variables as variables} \\
  k & && \tag{Natural numbers, $\N$} \\
  \delta &:=&& \cdot_{g?}^k \sep \wk_g^k \sep \delta, t/x \tag{Local substitutions} \\
  B &:=&&\ g : \Ctx \sep U : \DTyp i l \sep u : \DTrm i T l
          \tag{Global bindings} \\
  \Phi, \Psi &:= &&\ \cdot \sep \Phi, B \tag{Global contexts} \\
  S, T & := && \cdots \sep U^\delta \sep \CPI g l T \sep \TPI U l{l'}T \sep \CTyp l \sep \CTrm T l \\
  s, t & := && \cdots \sep u^\delta \sep \CLAM l g t \sep \CAPP t \Gamma
               \sep \TLAM l{l'} U t \sep \TAPP t T
               \sep \boxit T \sep \boxit t \\
    & && \sep \LETBTYP{l'}l\Gamma{x_T. M}{U}{t'}t \sep \LETBTRM{l'}l\Gamma
         T{x_t. M}{u}{t'}t \\
    & && \sep \ELIMTYP{l_1}{l_2}Mbl\Gamma t \sep \ELIMTRM{l_1}{l_2}Mbl\Gamma T t \\
  \vect M & := &&\ (\ell,g,x_T. M_\Typ)~(\ell,g,U_T,x_t. M_\Trm)
                  \tag{Two motives for mutual recursion of code} \\
  \vect b &:= && \vect{b}_\Typ~\vect{b}_\Trm
                 \tag{Branches for mutual recursion of code} \\
  b_\Typ &:= &&\ (g. t_\Nat) \sep (\ell,\ell',g, U_S,U_T,x_S,x_T. t_\Pi) \sep (\ell,g. t_{\Se})
                \sep (\ell,g,u_t,x_t. t_\tEl)
                       \tag{Branches for code of types} \\
  b_\Trm &:= &&\ (\ell,g,U_T,u_x. t_x) \sep (g.t'_\Nat) \sep (\ell,\ell',g,u_s,u_t,x_s,x_t.t'_\Pi) \sep
                (\ell,g.t'_\Se) 
                \tag{Branches for code of terms} \\
    & && \sep (g.t_\ze) \sep (g,u_t,x_t.t_{\tsucc}) \sep
         (\ell,g,U_M,u_s,u_{s'},u_t,x_M,x_s,x_{s'},x_t.t_{\telimn}) \\
    & && \sep (\ell,\ell',g,U_S,U_T,u_t,x_S,x_t.t_\lambda) \sep (\ell,\ell',g,U_S,U_T,u_t,u_s,x_S,x_T,x_t,x_s.t_\tapp)
\end{alignat*}
Following the layering principle before, we index our judgments with an layer index
$i$. %
We include four layers and these layers are summarized in \Cref{tab:dt:feat}. %
To elaborate, we begin with layer $\varlevel$, which is the layer that contains only
variables. %
This layer is needed to describe the base case of the recursive principles for code
when a local variable is hit. %
Layer $\varlevel$ (for variables) is \emph{not} available for most rules given in
\Cref{sec:dt:mltt-rules} other than the local variable rule and its congruence. %
We will follow this convention for the rest of this report, unless we specifically
state that layer $\varlevel$ is available for particular rules. %
Layer $\codelevel$ (for code) is the layer for code of MLTT. %
This layer is akin to layer $0$ in \Cref{sec:cv,sec:cv:logrel}, where static code
resides and no computation is allowed. %
However, in order to capture code promotions, we must introduce another layer, $\proglevel$ (for
programs), to permit computation of MLTT \emph{programs} in local contexts and on the
type level. %
This layer is especially crucial in the recursive principle for terms for the argument
$T$ where code promotions are implicitly handled by definitional equivalence. %
However, no meta-programs are allowed at layer $\proglevel$; in other words, the language at
layer $\proglevel$ is virtually vanilla MLTT. %
Therefore, unlike simple types, there are two layers in \delamlang permitting
computation. %
At last, we have layer $\metalevel$ (for meta-programs) where we have the power to do
meta-programming. %
At this layer, we have access to not only universe-polymorphic functions, but also
contextual types and recursive principles for code. %
All meta-functions must reside at this layer. %
All layers are related by a strict order of $\varlevel < \codelevel < \proglevel < \metalevel$. %

The reason to introduce layer $\proglevel$ becomes obvious by considering which layer the type
of a given MLTT term should live in. %
For instance, given a judgment $\lttyping \codelevel t T l$, we know $t$ lives at layer $\codelevel$ as
code, but what about $T$? %
Since $T$ is a type and we want $T$ to compute, it cannot live at layer $\codelevel$, but also
not $\metalevel$ as it must be a well-formed pure MLTT type. %
Indeed, $T$ ought to live at layer $\proglevel$, i.e. $\lttypwf \proglevel T l$. %
What about $\lttyping \proglevel t T l$? %
In this case, $t$ lives at layer $\proglevel$. %
Since $T$ must still be a well-formed pure MLTT type and compute, we must have
$\lttypwf \proglevel T l$. %
The type of a term living at layer $\metalevel$ simply also lives at layer $\metalevel$.  %
The relation of layers of terms and types is summarized by the following function:
\begin{align*}
  \typeof \varlevel &:= \proglevel \\
  \typeof \codelevel &:= \proglevel \\
  \typeof \proglevel &:= \proglevel \\
  \typeof \metalevel &:= \metalevel
\end{align*}
The judgment $\compt i$ quantifies computable layers:
\begin{mathpar}
  \inferrule
  { }
  {\compt \proglevel}

  \inferrule
  { }
  {\compt \metalevel}
\end{mathpar}

Then we extend our system with five types, following~\Cref{sec:cv}:
\begin{itemize}
\item $U^\delta$ is a global variable for types. %
  Due to separation of types and terms, we need a way to refer to code of types on the
  type level. 
\item $\CPI g l T$ is a meta-function type for introducing a context variable $g$
  to the global context. %
  We also have this in~\Cref{sec:cv}.
\item $\TPI U l{l'}T$ is similarly a meta-function type for introducing a type at
  layer $\proglevel$ to the global context. %
  This type is introduced to provide an index for the contextual type for terms to be
  discussed in the second next item.
\item $\CTyp l$ is a contextual type for types in MLTT. %
  It represents a static code of types. %
\item Finally, $\CTrm T l$ is a contextual type for terms in MLTT. %
  This $T$ may refer to the index type at layer $\proglevel$ introduced by meta-functions above. 
\end{itemize}

Since there are two kinds of contextual types now, there four kinds of bindings in a
global context:
\begin{itemize}
\item context variables $g : \Ctx$ representing a local context;
\item global variables $U : \DTyp i l$ representing a type in MLTT (note that $i \in
  \{\codelevel, \proglevel\}$);
\item global variables $u : \DTrm i T l$ representing a term in MLTT (note that
  $i \in \{\varlevel, \codelevel\}$ and there is no way to introduce a term at layer $\proglevel$ to global
  context).
\end{itemize}

Now, let us move on to discuss the extended terms. %
\begin{itemize}
\item First, we also introduce global variables and local substitutions. %
  Their syntax is identical to one in \Cref{sec:cv}. 
\item Then we have the introduction and elimination forms for meta-functions of
  context variables, $\CLAM l g t$ and $\CAPP t \Gamma$. 
\item Similarly, we have the introduction and elimination forms for meta-functions of
  types, $\TLAM l{l'} U t$ and $\TAPP t T$. 
\item Then we have the introduction forms of two kinds of contextual types.
\item Following \citet[Sec. 4]{hu2024layered}, we have two elimination forms for each kind of
  contextual types, $\tletbox$ and the recursive principles. %
  Same as before, $\tletbox$ is responsible for code composition and evaluation. %
  Intentional analyses are done through the recursive principles. %
  In \delamlang, $\tletbox$ is a bit more complex as it requires a specified motive. %
  We alter the syntax a little bit to make $\tletbox$ more like an operation:
  $\LETBTYP{l'}l\Gamma{x_T. M}{u}{t'}t$ and $\LETBTRM{l'}l\Gamma T{x_t. M}{u}{t'}t$.

\item Finally, we extend the recursive principles for code. %
  As indicated before, code of types and terms in MLTT are mutually defined, so the
  recursive principles must also be mutual. %
  The two recursive principles $\ELIMTYP{l_1}{l_2}Mbl\Gamma t$ and
  $\ELIMTRM{l_1}{l_2}Mbl\Gamma T t$ require two motives, one for types and one for terms,
  and contain all branches for code of types and terms. %
  Their difference is what exactly eventually being eliminated, as indicated by their
  subscript. %
  The branches $\vect b$ are a list of branches $\vect b_\Typ$ and $\vect b_\Trm$,
  where $\vect b_\Typ$ and $\vect b_\Trm$ contain \emph{all} branches for types and
  terms, respectively. %

\end{itemize}
In the branches, there are four kinds of variables.
\begin{itemize}
\item There is a globally introduced context variable $g$ which represents the
  local context where the code lives.
\item There could be some universe variables that are used to tell the universe levels
  of some types.
\item There could be some global variables $u$ and $U$ represent the sub-structures of
  a given case. %
  The subscripts correspond tightly to the syntax given in \Cref{sec:dt:syntax}.
\item For each sub-structure, there is one corresponding recursive variable $x$. %
  Again, the subscripts correspond tightly to the sub-structure. 
\end{itemize}

\subsection{More Typing and Equivalence Judgments}

In this section, we specify the remainder of the rules. %
We begin with the well-formedness rule for the global contexts. %
Recall that layer $\varlevel$ does not apply for most rules below, unless the otherwise is
specifically stated.
\begin{mathpar}
  \inferrule
  { }
  {\judge[L]\cdot}

  \inferrule
  {\judge[L]\Psi}
  {\judge[L]{\Psi, g : \Ctx}}

  \inferrule
  {\judge[L]\Psi \\ \lpjudge \proglevel \Gamma \\\\ \typing[L]l\Level \\ i \in \{\codelevel, \proglevel\}}
  {\judge[L]{\Psi, U : \DTyp i l}}

  \inferrule
  {\judge[L]\Psi \\ \lttypwf \proglevel T l \\\\ \typing[L]l\Level \\ i \in \{\varlevel, \codelevel\}}
  {\judge[L]{\Psi, u : \DTrm i T l}}
\end{mathpar}
The judgments for local substitutions follow very closely \Cref{sec:cv}. %
In these rules, $i$ might take $\varlevel$. %
This permission has a particular effect on the step case, which forces all terms in a
local substitution must be variables. %
\begin{mathpar}
  \inferrule
  {\lpjudge{\typeof i}\Gamma \\\\ \text{$\Gamma$ ends with $\cdot$} \\ |\Gamma| = k}
  {\ltsubst i {\cdot^k}{\cdot} \\ \ltsubeq i {\cdot^k}{\cdot^k}{\cdot}}

  \inferrule
  {\lpjudge{\typeof i}\Gamma \\ g : \Ctx \in \Psi \\\\ \text{$\Gamma$ ends with $g$} \\ |\Gamma| = k}
  {\ltsubst i {\cdot_g^k}{\cdot} \\ \ltsubeq i {\cdot_g^k}{\cdot_g^k}{\cdot}}

  \inferrule
  {\lpjudge{\typeof i}\Gamma \\ g : \Ctx \in \Psi \\ \text{$\Gamma$ ends with $g$} \\ |\Gamma| = k}
  {\ltsubst i {\wk_g^k}{g} \\ \ltsubeq i {\wk_g^k}{\wk_g^k}{g}}

  \inferrule
  {\typing[L]l\Level \\ \ltsubst i {\delta}{\Delta} \\ \lttypwf[\Psi][\Delta]{\typeof i}T l \\ \lttyping i {t}{T[\delta]} l}
  {\ltsubst i {\delta, t/x}{\Delta, x : T \at l}}

  \inferrule
  {\typing[L]l\Level \\ \ltsubeq i {\delta}{\delta'}{\Delta} \\ \lttypwf[\Psi][\Delta]{\typeof i}T l \\
    \lttyping i {t}{T[\delta]} l \\ \lttyping i {t'}{T[\delta]} l \\ \lttyequiv i {t}{t'}{T[\delta]} l}
  {\ltsubeq i {\delta, t/x}{\delta', t'/x}{\Delta, x : T \at l}}
\end{mathpar}
In the step case for equivalence above, we ask for two redundant premises of the
well-typedness of $t$ and $t'$ to provide an early presupposition for equivalence of
local substitutions. %
We will need this early presupposition in \Cref{lem:dt:lsubst-eq}. %
It breaks the dependency loop so that reaching the full presupposition lemma becomes
viable. 

Now let us consider the extended types and their equivalence.
\begin{mathpar}
  \inferrule
  {\lpjudge{\typeof i}\Gamma \\ U : \DTyp[\Delta]{i'} l \in \Psi \\ i' \in \{\codelevel, \proglevel\} \\ i' \le i \\ \ltsubst i \delta \Delta}
  {\lttypwf i{U^\delta}{l}}

  \inferrule
  {\lpjudge \metalevel\Gamma \\\\ \lttypwf[\Psi, g : \Ctx] \metalevel T l \\ \typing[L]l\Level}
  {\lttypwf \metalevel{\CPI g l T}{l}}

  \inferrule
  {\lpjudge \metalevel\Gamma \\ \lttypwf[\Psi, U : \DTyp[\Delta] \proglevel l] \metalevel T{l'} \\\\ \typing[L]l\Level \\ \typing[L]{l'}\Level}
  {\lttypwf \metalevel{\TPI U[\Delta] l{l'} T}{l'}}

  \inferrule
  {\lpjudge \metalevel\Gamma \\ \lpjudge \proglevel \Delta \\ \typing[L]l\Level}
  {\lttypwf \metalevel{\CTyp[\Delta]l}{0}}

  \inferrule
  {\lpjudge \metalevel\Gamma \\ \lttypwf[\Psi][\Delta] \proglevel T l \\ \typing[L]l\Level}
  {\lttypwf \metalevel{\CTrm[\Delta]T l}{0}}
\end{mathpar}
Notice that contextual types always live on the universe level $0$. %
This is fine because the meta-language at layer $\metalevel$ is strictly more expressive than
that at layer $\codelevel$, i.e. \mltt. %
Therefore, the universe level $0$ is large enough to encode all intrinsically typed
syntax of \mltt without causing any size issue. %
This intuition is further justified in the logical relations, where the semantics of
contextual types need not to refer to other semantics at layer $\metalevel$, so contextual
types can safely reside on level $0$. %
Similarly, meta-functions with open types as inputs also do not need to be concerned
about the universe levels of the types. %

The additional equivalence rules are just their congruence rules:
\begin{mathpar}
  \inferrule
  {\lpjudge{\typeof i}\Gamma \\ U : \DTyp[\Delta]{i'} l \in \Psi \\ i' \in \{\codelevel, \proglevel\} \\ i' \le i \\ \ltsubeq i \delta{\delta'} \Delta}
  {\lttypeq i{U^\delta}{U^{\delta'}}{l}}

  \inferrule
  {\lpjudge \metalevel\Gamma \\ \lttypeq[\Psi, g : \Ctx] \metalevel T{T'} l \\ \tyequiv[L]l{l'}\Level}
  {\lttypeq \metalevel{\CPI g l T}{\CPI g{l'}{T'}}{l}}

  \inferrule
  {\lpjudge \metalevel\Gamma \\ \lpequiv \proglevel \Delta{\Delta'} \\ \lttypeq[\Psi, U : \DTyp[\Delta] \proglevel l] \metalevel T{T'}{l'} \\ \tyequiv[L]{l_1}{l_3}\Level \\ \tyequiv[L]{l_2}{l_4}\Level}
  {\lttypeq \metalevel{\TPI U[\Delta]{l_1}{l_2}T}{\TPI U[\Delta']{l_3}{l_4}{T'}}{l_2}}

  \inferrule
  {\lpjudge \metalevel\Gamma \\ \lpequiv \proglevel \Delta{\Delta'} \\ \tyequiv[L]l{l'}\Level}
  {\lttypeq \metalevel{\CTyp[\Delta]l}{\CTyp[\Delta']{l'}}{0}}

  \inferrule
  {\lpjudge \metalevel\Gamma \\ \lpequiv \proglevel \Delta{\Delta'} \\ \lttypeq[\Psi][\Delta] \proglevel T{T'} l \\ \tyequiv[L]l{l'}\Level}
  {\lttypeq \metalevel{\CTrm[\Delta]T l}{\CTrm[\Delta']{T'}{l'}}{0}}
\end{mathpar}

Next, we list the extended typing judgments:
\begin{mathpar}
  \inferrule
  {\lpjudge{\typeof i}\Gamma \\ u : \DTrm[\Delta]{i'}T l \in \Psi \\ i' \in \{\varlevel, \codelevel\} \\ i \in
    \{\varlevel, \codelevel, \proglevel, \metalevel\} \\ i' \le i \\ \ltsubst i \delta \Delta}
  {\lttyping i{u^\delta}{T[\delta]}{l}}

  \inferrule
  {\lpjudge \metalevel\Gamma \\\\ \lttyping[\Psi, g : \Ctx]\metalevel t{T}{l} \\ \typing[L]l\Level}
  {\lttyping \metalevel{\CLAM l g t}{\CPI g l T}{l}}

  \inferrule
  {\lttypwf[\Psi, g : \Ctx] \metalevel T l  \\\\ \lttyping \metalevel{t}{\CPI g l T}{l} \\ \lpjudge \proglevel \Delta}
  {\lttyping \metalevel{\CAPP t \Delta}{T[\Delta/g]}{l}}

  \inferrule
  {\lpjudge \metalevel\Gamma \\ \lttyping[\Psi, U : \DTyp[\Delta] \proglevel l]\metalevel{t}{T}{l'} \\ \typing[L]{l}\Level \\ \typing[L]{l'}\Level}
  {\lttyping \metalevel{\TLAM{l}{l'}U t}{\TPI U[\Delta] l{l'} T}{l'}}

  \inferrule
  {\lttypwf[\Psi, U : \DTyp[\Delta] \proglevel l] \metalevel {T'}{l'} \\ \lttyping \metalevel{t}{\TPI U[\Delta] l{l'}{T'}}{l'} \\ \lttypwf[\Psi][\Delta] \proglevel{T} l}
  {\lttyping \metalevel{\TAPP t{T}}{T'[T/U]}{l'}}

  \inferrule
  {\lpjudge \metalevel \Gamma \\ \lttypwf[\Psi][\Delta]\codelevel T l}
  {\lttyping \metalevel{\boxit T}{\CTyp[\Delta] l}{0}}

  \inferrule
  {\lpjudge \metalevel \Gamma \\ \lttyping[\Psi][\Delta]\codelevel t T l}
  {\lttyping \metalevel{\boxit t}{\CTrm[\Delta] T l}{0}}

  \inferrule
  {\typing[L]{l'}\Level \\ \typing[L]l\Level \\ \lpjudge \proglevel \Delta \\
    \lttyping \metalevel{t}{\CTyp[\Delta]l}{1 + l} \\
    \lttypwf[\Psi][\Gamma,x_T : \CTyp[\Delta]l \at{0}]\metalevel{M}{l'} \\
  \lttyping[\Psi, U : \DTyp[\Delta]\codelevel l]\metalevel{t'}{M[\boxit U/x_T]}{l'}}
  {\lttyping \metalevel{\LETBTYP{l'}l \Delta{x_T.M}{U}{t'}t}{M[t/x_T]}{l'}}

  \inferrule
  {\typing[L]{l'}\Level \\ \typing[L]l\Level \\ \lpjudge \proglevel \Delta \\
    \lttypwf[\Psi][\Delta]\proglevel T l \\
    \lttyping \metalevel{t}{\CTrm[\Delta] T l}{l} \\
    \lttypwf[\Psi][\Gamma,x_t : \CTrm[\Delta] T l \at{0}]\metalevel{M}{l'} \\
  \lttyping[\Psi, u : \DTrm[\Delta]\codelevel T l]\metalevel{t'}{M[\boxit u/x_t]}{l'}}
  {\lttyping \metalevel{\LETBTRM{l'}l \Delta T{x_t.M}{u}{t'}t}{M[t/x_t]}{l'}}
\end{mathpar}
For the typing rule of $\CAPP t \Delta$, we require $\Delta$ to be a context at layer
$\proglevel$. %
This is because a context variable represents a local context for an MLTT term. %
A local context for an MLTT term necessarily lives at layer $\proglevel$, so we can only
substitute a context living at layer $\proglevel$ with a context variable. %

Now we shall mentally prepare ourselves to write down the typing rules for the two
recursive principles. %
They are conceptually easy but simply verbose to write down. %
We will only write down the rules for this time for completeness and in later
discussions, we simply omit the premises. %
% As both recursive principles share most premises, we write them
Our goal is to provide the following conclusions:
\begin{gather*}
  \lttyping \metalevel {\ELIMTYP{l_1}{l_2}M b{l'}\Delta t}{M_\Typ[l'/l,\Delta/g,t/x_T]}{l_1} \\
  \lttyping \metalevel {\ELIMTRM{l_1}{l_2}M b{l'}\Delta T t}{M_\Trm[l'/l,\Delta/g,T/U_T,t/x_t]}{l_2}
\end{gather*}
We group the premises into different parts. %
First we give the premises related to the motives:
\begin{mathpar}
  \typing[L]{l_1}\Level

  \typing[L]{l_2}\Level

  \lttypwf[\Psi, g : \Ctx][\Gamma, x_T : \CTyp[g]\ell \at{0}][L,\ell] \metalevel{M_\Typ}{l_1}

  \lttypwf[\Psi, g : \Ctx, U_T : \DTyp[g]\proglevel\ell][\Gamma, x_t : \CTrm[g]{U_T^\id}\ell \at{0}][L,\ell] \metalevel{M_\Trm}{l_2}
\end{mathpar}
where $\vect M =  (\ell,g,x_T. M_\Typ)~(\ell,g,U_T,x_t. M_\Trm)$. %
In the premises above, we give the well-formedness of two motives for code of types
and terms, respectively. %
Let us call this group $G_M$. %
We move on to considering the branches. %
We first consider the branches for code of types. %
It is relatively easy as there are only four cases:
\begin{itemize}
\item 
  \[
    \lttyping[\Psi, g : \Ctx]\metalevel{t_\Nat}{M_\Typ[0/\ell,g/g,\boxit \Nat/x_T]}{l_1}
  \]
\item We explain this premise more carefully. Consider some code of type
  $\boxit{\PI{l}{l'}x S T}$, then we have the matching premise
  \[
    \lttyping[\Psi'][\Gamma'][L,\ell,\ell']\metalevel{t_\Pi}{M_\Typ[\ell\sqcup \ell'/\ell,g/g,\boxit{\PI{\ell}{\ell'}x{U_S^\id}{U_T^\id}}/x_T]}{l_1}
  \]
  where
  \begin{align*}
    \Psi' :=&~ \Psi \\
    ,&~ g : \Ctx \\
    ,&~U_S : \DTyp[g] \codelevel \ell
       \tag{the global variable for the input type, which captures $S$}\\
    ,&~U_T :\DTyp[g,x:U_S^\id \at{\ell}]\codelevel {\ell'}
       \tag{the global variable for the output type, which captures $T$; note that it
       lives in an extended local context}
  \end{align*}
  and
  \begin{align*}
    \Gamma' :=&~ \Gamma \\
    ,&~ x_S: M_\Typ[\ell/\ell,g/g,\boxit{U_S^\id}/x_T] \at{l_1}
       \tag{the recursive call for $S$ of type $M$ that is properly substituted}\\
    ,&~x_T:M_\Typ[\ell'/\ell,(g,x:U_S^\id \at{\ell})/g,\boxit{U_T^\id}/x_T]\at {l_1}
       \tag{the recursive call for $T$; see how the local context is extended}
  \end{align*}
\item
  Further,
  \[
    \lttyping[\Psi,g:\Ctx][\Gamma][L,\ell]\metalevel{t_\Se}{M_\Typ[1 + \ell/\ell,g/g,\boxit{\Ty \ell}/x_T]}{l_1}
  \]
\item
  \[
    \lttyping[\Psi,g:\Ctx, u_t:\DTrm[g]\codelevel{\Ty \ell}{1 + \ell}][\Gamma'][L,\ell]\metalevel{t_\tEl}{M_\Typ[\ell/\ell,g/g,\boxit{(\Elt \ell{u_t^\id})}/x_T]}{l_1}
  \]
  where
  \[
    \Gamma' := \Gamma, x_t : M_\Trm[1 + \ell/\ell,g/g,\Ty \ell/U_T,\boxit{u_t^\id}/x_t]  \at{l_2}
  \]
\end{itemize}
Let us call this group $G_\Typ$. 

Lastly, let us consider the nine cases for terms.
\begin{itemize}
\item
  \[
    \lttyping[\Psi, g : \Ctx,U_T : \DTyp[g] \proglevel \ell, u_x : \DTrm[g]\varlevel{U_T^\id}{\ell}][\Gamma][L,\ell]\metalevel{t_x}{M_\Trm[\ell/\ell,g/g,U_T^\id/U_T,\boxit u_x/x_t]}{l_2}
  \]
  In this case, the type of the variable is captured by $U_T$. %
  It has to live at layer $\proglevel$ because it is not a sub-structure of the variable,
  i.e. it is obtained externally, from the indexing arguments of the recursive principles. %
\item
  \[
    \lttyping[\Psi, g : \Ctx]\metalevel{t'_\Nat}{M_\Trm[1/\ell,g/g,\Ty 0/U_T,\boxit \Nat/x_t]}{l_2}    
  \]
\item
  \[
    \lttyping[\Psi'][\Gamma'][L,\ell,\ell']\metalevel{t'_\Pi}{M_\Trm[1 +{(\ell\sqcup \ell')}/\ell,g/g,\Ty{\ell\sqcup \ell'}/U_T,\boxit{\PI{\ell}{\ell'}x{u_s^\id}{u_t^\id}}/x_t]}{l_2}
  \]
  where
  \[
    \Psi' := \Psi, g : \Ctx,u_s : \DTrm[g] \codelevel {\Ty \ell}{1 + \ell},u_t :\DTrm[g,x:\Elt{\ell}{u_s^\id}]\codelevel{\Ty{\ell'}} {1 + \ell'}
  \]
  and
  \begin{align*}
    \Gamma' :=&~ \Gamma \\
    ,&~ x_s: M_\Trm[1 + \ell/\ell,g/g,\Ty \ell/U_T,\boxit{u_s^\id}/x_t] \at{l_2} \\
    ,&~ x_t:M_\Trm[1 + \ell'/\ell,(g,x:\Elt{\ell}{u_s^\id})/g,\Ty{\ell'}/U_T,\boxit{u_t^\id}/x_t]\at {l_2}
  \end{align*}
  Notice that this premise for the encoding of $\Pi$ is almost identical to the
  premise in $G_\Typ$ above, with necessary adjustment to return the proper motive
  $M_\Trm$ instead. %
  
\item
  \[
    \lttyping[\Psi,g:\Ctx][\Gamma][L,\ell]\metalevel{t'_\Se}{M_\Trm[{2 + \ell}/\ell,g/g, \Ty{1 + \ell}/U_T,\boxit{\Ty \ell}/x_T]}{l_2}
  \]
\item
  \[
    \lttyping[\Psi, g : \Ctx]\metalevel{t_\ze}{M_\Trm[0/\ell,g/g,\Nat/U_T,\boxit \ze/x_t]}{l_2}    
  \]
\item 
  \[
    \lttyping[\Psi, g : \Ctx,u_t: \DTrm[g]\codelevel\Nat 0][\Gamma']\metalevel{t_\tsucc}{M_\Trm[0/\ell,g/g,\Nat/U_T,\boxit{(\su{u_t^\id})}/x_t]}{l_2}    
  \]
  where $\Gamma' := \Gamma, x_t:M_\Trm[0/\ell,g/g,\Nat/U_T,\boxit{u_t^\id}/x_t] \at{l_2}$.
\item We carefully explain this premise for the code of the elimination of natural
  numbers. %
  Recall that the syntax is $\ELIMN l{x.M}{s}{x,y.s'}t$. %
  We use corresponding global variables to capture the sub-structures. 
  \[
    \lttyping[\Psi'][\Gamma'][L,\ell]\metalevel{t_{\telimn}}{M_\Trm[\ell/\ell,g/g,U_M^{\id_g,u_t^\id/x}/U_T,\ELIMN \ell{x.U_M^{\id_{g,x}}}{u_s^{\id_g}}{x,y. u_{s'}^{\id_{g,x,y}}}{u_t^{\id_g}}/x_t]}{l_2}
  \]
  where
  \begin{align*}
    \Psi' :=&~ \Psi \\
    ,&~g : \Ctx \\
    ,&~U_M : \DTyp[g,x:\Nat\at 0]\codelevel \ell
       \tag{the global variable for the motive; it lives at layer $\codelevel$ as it is a sub-structure}\\
    ,&~u_s : \DTrm[g]\codelevel{U_M^{\id_g,\ze/x}}\ell
       \tag{the code for the base case; its type refers to the code of the motive with
       $x$ for $\ze$}\\
    ,&~u_{s'} : \DTrm[g, x : \Nat\at 0, y : U_M^{\id_{g,x}} \at
       \ell]\codelevel{U_M^{\id_g,\su{x}/x}}\ell
       \tag{the code for the step case; the local context is extended with the
       predecessor and the recursive call} \\
    ,&~u_t : \DTrm[g]\codelevel{\Nat}0
       \tag{the code for the scrutinee}
  \end{align*}
  and
  \begin{align*}
    \Gamma' :=&~\Gamma \\
    ,&~x_M : M_\Typ[\ell/\ell,(g,x:\Nat\at 0)/g,\boxit{U_M^{\id_{g,x}}}/x_T] \at{l_1}
       \tag{since the motive is a sub-structure, a recursive call is available}\\
    ,&~x_s : M_\Trm[\ell/\ell,g/g,U_M^{\id_g,\ze/x}/U_T,\boxit{u_s^{\id_g}}/x_t] \at{l_2}
       \tag{the recursive call for the base case; recall that $U_T$ is the type of
       $x_t$,}\\
    \tag{which in this case is also the type of $u_s$} \\
    ,&~x_{s'} : M_\Trm[\ell/\ell,(g, x : \Nat\at 0, y : U_M^{\id_{g,x}} \at
       \ell)/g,U_M^{\id_g,\su x/x}/U_T,\boxit{u_{s'}^{\id_{g,x,y}}}/x_t] \at{l_2}
       \tag{the recursive call for the step case; similar logic applies but more longer}\\
    ,&~x_t : M_\Trm[\ell/\ell,g/g,\Nat/U_T,\boxit{u_t^\id}/x_t] \at{l_2}
       \tag{the recursive call for the scrutinee}
  \end{align*}
\item
  \[
    \lttyping[\Psi'][\Gamma'][L,\ell,\ell']\metalevel{t_\lambda}{M_\Trm[\ell\sqcup \ell'/\ell,g/g,\PI
      \ell{\ell'}x{U_S^{\id_g}}{U_T^{\id_{g,x}}}/U_T, \boxit{\LAM \ell{\ell'} x{U_S^{\id_g}}{u_t^{\id_{g,x}}}}/x_t]}{l_2}
  \]
  where
  \begin{align*}
    \Psi' :=&~ \Psi \\
    ,&~g : \Ctx \\
    ,&~U_S : \DTyp[g]\codelevel \ell \\
    ,&~U_T : \DTyp[g, x : U_S^{\id_g} \at \ell]\proglevel{\ell'} \\
    ,&~u_t : \DTrm[g, x : U_S^{\id_g} \at \ell]\codelevel{U_T^{\id_{g,x}}}{\ell'}
  \end{align*}
  and
  \begin{align*}
    \Gamma' :=&~\Gamma \\
    ,&~x_S : M_\Typ[\ell/\ell,g/g,\boxit{U_S^{\id_g}}/x_T] \at{l_1} \\
    ,&~x_t : M_\Trm[\ell'/\ell,(g,x : U_S^{\id_g} \at \ell)/g,U_T^{\id_{g,x}}/U_T,\boxit{u_t^{\id_{g,x}}}/x_t] \at{l_2}
  \end{align*}
  Note that here $U_T$ is at layer $\proglevel$. %
  This is because the return type of not a sub-structure in a function abstraction
  $\LAM l{l'} x S t$, and therefore it must be captured externally from the indexing
  arguments of the recursive principle. %
  Since it is not a sub-structure, there also is not a recursive call for it. %
  It is possible to include the return type as a sub-structure, e.g. $\LAM{l}{l'}x S{(t
    : T)}$ but we decided to show this alternative to demonstrate various design spaces.
  
\item Finally,
  \[
    \lttyping[\Psi'][\Gamma'][L,\ell,\ell']\metalevel{t_\tapp}{M_\Trm[\ell'/\ell,g/g,U_T^{\id_g,u_t^{\id_g}/x}/U_T,\boxit{(\APP{u_t^{\id_g}}
        \ell{\ell'} x{U_S^{\id_{g}}}{U_T^{\id_{g,x}}}{u_s^{\id_g}})}/x_t]}{l_2}
  \]
  where
  \begin{align*}
    \Psi' :=&~ \Psi \\
    ,&~g : \Ctx \\
    ,&~U_S : \DTyp[g]\codelevel \ell \\
    ,&~U_T : \DTyp[g, x : U_S^{\id_g} \at \ell]\codelevel{\ell'} \\
    ,&~u_t : \DTrm[g]\codelevel{\PI \ell{\ell'}x{U_S^{\id_{g}}}{U_T^{\id_{g,x}}}}{\ell \sqcup \ell'} \\
    ,&~u_s : \DTrm[g]\codelevel{U_S^{\id_{g}}}{\ell}
  \end{align*}
  and
  \begin{align*}
    \Gamma' :=&~\Gamma \\
    ,&~x_S : M_\Typ[\ell/\ell,g/g,\boxit{U_S^{\id_g}}/x_T] \at{l_1} \\
    ,&~x_T : M_\Typ[\ell/\ell,(g,x : U_S^{\id_g} \at \ell)/g,\boxit{U_T^{\id_{g,x}}}/x_T] \at{l_1} \\
    ,&~x_t : M_\Trm[\ell \sqcup \ell'/\ell,g/g,\PI \ell{\ell'}x{U_S^{\id_{g}}}{U_T^{\id_{g,x}}}/U_T,\boxit{u_t^{\id_{g}}}/x_t] \at{l_2} \\
    ,&~x_s : M_\Trm[\ell/\ell,g/g,U_S^{\id_{g}}/U_T,\boxit{u_s^{\id_{g}}}/x_t] \at{l_2}
  \end{align*}
  This premise shows why we must use a more verbose syntax for application, i.e. $\APP
  t l{l'}xSTs$. %
  In the global context, we must introduce the global variables for the input and
  output types. %
  However, a vanilla function application $t~s$ has no such information at all. %
  Since the current syntax has both input and output types as sub-structures, we can
  also allow their recursive calls. 
\end{itemize}
All premises above conclude the group for terms, which we name $G_\Trm$. %
We collectively use $G_A$ for all three groups above, i.e. $G_A :=
G_M~G_\Typ~G_\Trm$. %
Then we have the typing rule for the recursive principles as follows:
\begin{mathpar}
  \inferrule
  {G_A \\ \typing[L]{l'}\Level \\ \lpjudge \proglevel \Delta \\ \lttyping \metalevel t{\CTyp[\Delta]{l'}}{0}}
  {\lttyping \metalevel {\ELIMTYP{l_1}{l_2}M b{l'}\Delta t}{M_\Typ[l'/\ell,\Delta/g,t/x_T]}{l_1}}

  \inferrule
  {G_A \\ \typing[L]{l'}\Level \\ \lpjudge \proglevel \Delta \\ \lttypwf[\Psi][\Delta]\proglevel T{l'} \\
    \lttyping \metalevel t{\CTrm[\Delta]T{l'}}{0}}
  {\lttyping \metalevel {\ELIMTRM{l_1}{l_2}M b{l'}\Delta T t}{M_\Trm[l'/\ell,\Delta/g,T/U_T,t/x_t]}{l_2}}
\end{mathpar}


\subsection{More Congruence Rules for Typing}

The congruence rules for the additional typing rules are naturally derived from
the typing rules above.
\begin{mathpar}
  \inferrule
  {\lpjudge{\typeof i}\Gamma \\ u : \DTrm[\Delta]{i'}T l \in \Psi \\ i' \in \{\varlevel, \codelevel\} \\ i \in
    \{\varlevel, \codelevel, \proglevel, \metalevel\} \\ i' \le i \\ \ltsubeq i \delta{\delta'} \Delta}
  {\lttyequiv i{u^\delta}{u^{\delta'}}{T[\delta]}{l}}

  \inferrule
  {\lpjudge \metalevel\Gamma \\ \lttyequiv[\Psi, g : \Ctx]\metalevel t{t'}{T}{l} \\ \tyequiv[L]l{l'}\Level}
  {\lttyequiv \metalevel{\CLAM l g t}{\CLAM{l'} g{t'}}{\CPI g l T}{l}}

  \inferrule
  {\lttyequiv \metalevel{t}{t'}{\CPI g l T}{l} \\ \lpequiv \proglevel \Delta{\Delta'}}
  {\lttyequiv \metalevel{\CAPP t \Delta}{\CAPP{t'}{\Delta'}}{T[\Delta/g]}{l}}

  \inferrule
  {\lpjudge \metalevel\Gamma \\ \lttyequiv[\Psi, U : \DTyp[\Delta] \proglevel{l_1}]\metalevel{t}{t'}{T}{l_2} \\ \tyequiv[L]{l_1}{l_3}\Level \\ \tyequiv[L]{l_2}{l_4}\Level}
  {\lttyequiv \metalevel{\TLAM{l_1}{l_2}U t}{\TLAM{l_3}{l_4}U{t'}}{\TPI U[\Delta]{l_1}{l_2}
      T}{l_2}}

  \inferrule
  {\lttyequiv \metalevel{t}{t'}{\TPI U[\Delta] l{l'}{T''}}{l'} \\ \lttypeq[\Psi][\Delta] \proglevel{T}{T'}l}
  {\lttyequiv \metalevel{\TAPP t{T}}{\TAPP{t'}{T'}}{T''[T/U]}{l'}}
\end{mathpar}

The following rules are related to meta-programming and intensional analysis. 
\begin{mathpar}
  \inferrule
  {\lpjudge \metalevel \Gamma \\ \lttypeq[\Psi][\Delta]\codelevel T{T'}l}
  {\lttyequiv \metalevel{\boxit T}{\boxit{T'}}{\CTyp[\Delta] l}{0}}

  \inferrule
  {\lpjudge \metalevel \Gamma \\ \lttyequiv[\Psi][\Delta]\codelevel t{t'} T l}
  {\lttyequiv \metalevel{\boxit t}{\boxit{t'}}{\CTrm[\Delta] T l}{0}}

  \inferrule
  {\lpjudge \metalevel \Gamma \\ \tyequiv[L]{l_1}{l_3}\Level \\ \tyequiv[L]{l_2}{l_4}\Level \\ \lpequiv \proglevel \Delta{\Delta'} \\
    \lttyequiv \metalevel{t}{t'}{\CTyp[\Delta]{l_2}}{0} \\
    \lttypeq[\Psi][\Gamma,x_T : \CTyp[\Delta]{l_2} \at{0}]\metalevel{M}{M'}{l_1} \\
  \lttyequiv[\Psi, U : \DTyp[\Delta]\codelevel {l_2}]\metalevel{t_1}{t_2}{M[\boxit{U^\id}/x_T]}{l_1}}
  {\lttyequiv \metalevel{\LETBTYP{l_1}{l_2} \Delta{x_T.M}{U}{t_1}t}{\LETBTYP{l_3}{l_4}{\Delta'}{x_T.M'}{U}{t_2}{t'}}{M[t/x_T]}{l_1}}

  \inferrule
  {\lpjudge \metalevel \Gamma \\ \tyequiv[L]{l_1}{l_3}\Level \\ \tyequiv[L]{l_2}{l_4}\Level \\
    \lpequiv \proglevel \Delta{\Delta'} \\
    \lttypeq \proglevel T{T'}{l_2} \\
    \lttyequiv \metalevel{t}{t'}{\CTrm[\Delta]T{l_2}}{0} \\
    \lttypeq[\Psi][\Gamma,x_T : \CTrm[\Delta]T{l_2} \at{0}]\metalevel{M}{M'}{l_1} \\
    \lttyequiv[\Psi, u : \DTrm[\Delta]\codelevel T{l_2}]\metalevel{t_1}{t_2}{M[\boxit{u^\id}/x_t]}{l_1}}
  {\lttyequiv \metalevel{\LETBTRM{l_1}{l_2} \Delta{T}{x_t.M}{u}{t_1}t}{\LETBTRM{l_3}{l_4}{\Delta'}{T'}{x_T.M'}{u}{t_2}{t'}}{M[t/x_t]}{l_1}}
\end{mathpar}
We omit the congruence rules for the recursive principles for code as they are
conceptually simple but too long. %
We simply let equivalence to propagate inwards to all the sub-terms of the recursive
principles. 

\subsection{Computation Rules}

Finally, we list all the computation rules. %
In the rules below, we let $\compt i$. %
We first list the $\beta$ rules for
natural numbers:
\begin{mathpar}
  \inferrule
  {\typing[L]l\Level \\ \lttypwf[\Psi][\Gamma, x : \Nat \at 0] i M l \\
    \lttyping i s {M[\ze/x]}l \\
    \lttyping[\Psi][\Gamma, x : \Nat \at 0, y : M \at l] i {s'}{M[\su x/x]}l}
  {\lttyequiv i{s}{\ELIMN l{x.M}s{x,y. s'}\ze}{M[\ze/x]}{l}}

  \inferrule
  {\typing[L]l\Level \\ \lttypwf[\Psi][\Gamma, x : \Nat \at 0] i M l \\
    \lttyping i s {M[\ze/x]}l \\
    \lttyping[\Psi][\Gamma, x : \Nat \at 0, y : M \at l] i {s'}{M[\su x/x]}l \\
    \lttyping i t \Nat 0}
  {\lttyequiv i{s'[t/x,\ELIMN l{x.M}s{x,y. s'}t/y]}{\ELIMN l{x.M}s{x,y. s'}{(\su t)}}{M[\su t/x]}{l}}
\end{mathpar}
Then we have the $\beta$ and $\eta$ rules for dependent functions:
\begin{mathpar}
  \inferrule
  {\typing[L]l\Level \\ \typing[L]{l'}\Level \\ \lttypwf i S{l}
    \\
    \lttypwf[\Psi][\Gamma, x : S \at{l}] i{T}{l'} \\
    \lttyping[\Psi][\Gamma, x : S \at{l}] i t{T}{l'} \\ \lttyping i s{S}{l}}
  {\lttyequiv i{t[s/x]}{\APP {\LAM {l}{l'} x S t} {l}{l'} x{S} T s}{T[s/x]}{l'}}

  \inferrule
  {\typing[L]l\Level \\ \typing[L]{l'}\Level \\ \lttypwf i S{l} \\
    \lttypwf[\Psi][\Gamma, x : S \at{l}] i{T}{l'} \\
    \lttyping i t{\PI {l}{l'}x S T}{l \sqcup l'}}
  {\lttyequiv i{\LAM {l}{l'} x{S} {\APP t {l}{l'}x S T x}}{t}{\PI {l}{l'} x S T}{l \sqcup l'}}
\end{mathpar}
In the $\eta$ rule, on the right hand side, all $t$, $S$ and $T$ should be properly
locally weakened.

Finally we have $\beta$ and $\eta$ rules for universe-polymorphic functions:
\begin{mathpar}
  \inferrule
  {\lpjudge \metalevel \Gamma \\ \lttyping[\Psi][\Gamma][L, \vect \ell]\metalevel{t}{T}{l} \\ \typing[L,\vect\ell]l\Level \\
    |\vect\ell| = |\vect l| > 0 \\ \forall l' \in \vect l ~.~ \typing[L]{l'}\Level}
  {\lttyequiv \metalevel{t[\vect l/\vect \ell]}{(\ULAM l \ell t)~\$~\vect l}{T[\vect l/\vect \ell]}{l[\vect l/\vect \ell]}}

  \inferrule
  {\lttyping \metalevel{t}{\UPI \ell l T}{\omega}}
  {\lttyequiv \metalevel{\ULAM l \ell{(t~\$~\vect\ell)}}{t}{\UPI \ell l T}{\omega}}
\end{mathpar}
Similarly, in the $\eta$ rule, the universe variables appearing in $t$ must also be
properly weakened. 
This concludes all the rules for the MLTT portion of \delamlang. 

Then we move on to considering the computation rules for the extended types. %
Let us finish considering all meta-function types.
\begin{mathpar}
  \inferrule
  {\lpjudge \metalevel \Gamma \\ \lttyping[\Psi, g : \Ctx]\metalevel t{T}{l} \\ \typing[L]l\Level \\ \lpjudge \proglevel \Delta}
  {\lttyequiv \metalevel{t[\Delta/g]}{\CAPP{(\CLAM l g t)}\Delta}{T[\Delta/g]}{l}}

  \inferrule
  {\lttyping \metalevel{t}{\CPI g l T}{l}}
  {\lttyequiv \metalevel{\CLAM l g {(\CAPP t g)}}{t}{\CPI g l T}{l}}

  \inferrule
  {\lpjudge \metalevel \Gamma \\ \lttyping[\Psi, U : \DTyp[\Delta] \proglevel l]\metalevel{t}{T'}{l'} \\ \typing[L]{l}\Level \\ \typing[L]{l'}\Level \\ \lttypwf[\Psi][\Delta] \proglevel{T} l}
  {\lttyequiv \metalevel{t[T/U]}{\TAPP{(\TLAM{l}{l'}U t)}{T}}{T'[T/U]}{l'}}

  \inferrule
  {\lttyping \metalevel{t}{\TPI U[\Delta] l{l'}{T'}}{l'}}
  {\lttyequiv \metalevel{\TLAM{l}{l'}U{(\TAPP t{U^\id})}}{t}{\TPI U[\Delta] l{l'}{T'}}{l'}}
\end{mathpar}
Now we consider the contextual types. %
They only have $\beta$ rules. %
Let us consider $\tletbox$ first.
\begin{mathpar}
  \inferrule
  {\lpjudge \metalevel \Gamma \\ \typing[L]{l'}\Level \\ \typing[L]l\Level \\ \lpjudge \proglevel \Delta \\
    \lttypwf[\Psi][\Delta]\codelevel T l \\
    \lttypwf[\Psi][\Gamma,x_T : \CTyp[\Delta]l \at{0}]\metalevel{M}{l'} \\
  \lttyping[\Psi, U : \DTyp[\Delta]\codelevel l]\metalevel{t'}{M[\boxit U^\id/x_T]}{l'}}
  {\lttyequiv \metalevel{t'[T/U]}{\LETBTYP{l'}l \Delta{x_T.M}{U}{t'}{(\boxit T)}}{M[\boxit T/x_T]}{l'}}

  \inferrule
  {\lpjudge \metalevel \Gamma \\ \typing[L]{l'}\Level \\ \typing[L]l\Level \\ \lpjudge \proglevel \Delta \\
    \lttypwf[\Psi][\Delta]\proglevel T l \\
    \lttyping[\Psi][\Delta]\codelevel t T l \\
    \lttypwf[\Psi][\Gamma,x_t : \CTrm[\Delta] T l \at{0}]\metalevel{M}{l'} \\
  \lttyping[\Psi, u : \DTrm[\Delta]\codelevel T l]\metalevel{t'}{M[\boxit u^\id/x_t]}{l'}}
  {\lttyequiv \metalevel{t'[t/u]}{\LETBTRM{l'}l \Delta T{x_t.M}{u}{t'}{(\boxit t)}}{M[\boxit t/x_t]}{l'}}
\end{mathpar}
We can also give the $\beta$ rules for the recursive principles for code. %
There are too many to list them all, and moreover they follow the same pattern, so we
just list a selected few of them. %
We begin with something easy:
\begin{mathpar}
  \inferrule
  {G_A \\ \lpjudge \metalevel \Gamma \\ \lpjudge \proglevel \Delta}
  {\lttyequiv \metalevel {t_\Nat[\Delta/g]}{\ELIMTYP{l_1}{l_2}M b{0}\Delta{(\boxit
        \Nat)}}{M[0/\ell,\Delta/g,\boxit \Nat/x_T]}{l_1}}
\end{mathpar}
In this case, we provide $\Nat$ to the recursive principle for code of types. %
It hits the base case described by $t_\Nat$, and thus the whole program is reduced to
$t_\Nat$ with $g$ for $\Delta$. %
Note that the universe level for $\Nat$ must be $0$ as specified by the typing
judgment at layer $\codelevel$. %
The recursive principle for code of terms behaves very similarly when encountering the
code of $\Nat$. %
Instead, it picks the right branch $t'_\Nat$ and returns the right motive instead:
\begin{mathpar}
  \inferrule
  {G_A \\ \lpjudge \metalevel \Gamma \\ \lpjudge \proglevel \Delta}
  {\lttyequiv \metalevel {t'_\Nat[\Delta/g]}{\ELIMTRM{l_1}{l_2}M b{1}\Delta{\Ty 0}{(\boxit
        \Nat)}}{M_\Trm[1/\ell,\Delta/g,\Ty 0/U_T,\boxit \Nat/x_t]}{l_2}}
\end{mathpar}
In order to have the code to be $\boxit\Nat$ as a term, this code must have type $\Ty
0$, which lives at universe level $1$. %
Hence the indices are forced by the typing rules at layer $\codelevel$. 

Then we specify the variable case:
\begin{mathpar}
  \inferrule
  {G_A \\ \lpjudge \metalevel \Gamma \\ \typing[L]{l'}\Level \\ \lpjudge \proglevel \Delta \\ \lttypwf[\Psi][\Delta]\proglevel T{l'}
    \\
    x : T \at{l'} \in \Delta}
  {\lttyequiv \metalevel {t_x[l'/\ell,\Delta/g,T/U_T,x/u_x]}{\ELIMTRM{l_1}{l_2}M b{l'}\Delta{T}{(\boxit
        x)}}{M_\Trm[l'/\ell,\Delta/g,T/U_T,\boxit x/x_t]}{l_2}}
\end{mathpar}
The subtlety here is that $u_x$ can only receive a variable as it is typed at layer
$\varlevel$, but it is fine as $x$ is precisely just a variable. 

Then let us consider a more complex case of $\Pi$ types.
\begin{mathpar}
  \inferrule
  {G_A \\ \lpjudge \metalevel \Gamma \\ \typing[L]{l}\Level \\ \typing[L]{l'}\Level \\ \lpjudge \proglevel \Delta \\
    \lttypwf[\Psi][\Delta] \codelevel S l  \\ \lttypwf[\Psi][\Delta, x : S \at l]\codelevel{T}{l'} \\
    t = \boxit{\PI{l}{l'}x S T} \\
    s_S = \ELIMTYP{l_1}{l_2}M b{l}\Delta{(\boxit{S})} \\
    s_T = \ELIMTYP{l_1}{l_2}M b{l'}{(\Delta, x : S \at l)}{(\boxit{T})}}
  {\lttyequiv \metalevel {t_\Pi[l/\ell,l'/\ell',\Delta/g,S/U_S,T/U_T,s_S/x_S,
      s_T/x_T]}{\ELIMTYP{l_1}{l_2}M b{(l \sqcup
        l')}\Delta{t}}{M_\Trm[l\sqcup l'/\ell,\Delta/g,t/x_T]}{l_1}}
\end{mathpar}
Notice how $s_S$ and $s_T$ recurse down the sub-structures, i.e. $S$ and $T$ with the
proper universe levels and local contexts. %
We end our discussion by given the $\beta$ rules for code of function abstractions and
applications, as they appear to be rather complex, but their essence is fundamentally simple.
\begin{mathpar}
  \inferrule
  {G_A \\ \lpjudge \metalevel \Gamma \\ \lpjudge \proglevel \Delta \\ \typing[L]l\Level  \\
    \typing[L]{l'}\Level \\ \lttypwf[\Psi][\Delta] \codelevel S l \\
    \lttyping[\Psi][\Delta, x : S \at l] \codelevel t{T}{l'} \\
    l_\Pi = l \sqcup l' \\
    T_\Pi = \PI{l}{l'}x{S}{T} \\
    t' = \boxit{\LAM{l}{l'}x{S}{t}} \\
    s_S = \ELIMTYP{l_1}{l_2}M b{l}\Delta{(\boxit{S})} \\
    s_t = \ELIMTRM{l_1}{l_2}M b{l'}{(\Delta, x : S \at l)}{T}{(\boxit{t})} \\
    \delta = s_S/x_S,s_t/x_t}
  {\lttyequiv \metalevel {t_\lambda[l/\ell,l'/\ell',\Delta/g,S/U_S,T/U_T,t/u_t,\delta]}{\ELIMTRM{l_1}{l_2}M b{l_\Pi}\Delta{T_\Pi}{t'}}{M_\Trm[l_\Pi/\ell,\Delta/g,T_\Pi/U_T,t'/x_t]}{l_2}}
\end{mathpar}
Similarly, the recursive principle for code of terms picks the right branch
($t_\lambda$) with variables properly substituted. %
Since $S$ is also a sub-structure, the recursive call $s_S$ invokes the recursive
principle for code of types instead, hence making the recursive principles mutually
defined.

Last, we give the case for function applications.
\begin{mathpar}
  \inferrule
  {G_A \\ \lpjudge \metalevel \Gamma \\ \lpjudge \proglevel \Delta \\ \typing[L]l\Level  \\ \typing[L]{l'}\Level
    \\ \lttypwf[\Psi][\Delta]\codelevel S l \\ \lttypwf[\Psi][\Delta, x : S \at l] \codelevel{T}{l'} \\
    \lttyping[\Psi][\Delta] \codelevel t{\PI l{l'} x S T}{l \sqcup l'} \\
    \lttyping[\Psi][\Delta] \codelevel s S l \\
    T_\tapp = T[s/x] \\
    t' = \boxit{(\APP t l{l'} x S T s)} \\
    s_S = \ELIMTYP{l_1}{l_2}M b{l}\Delta{(\boxit{S})} \\
    s_T = \ELIMTYP{l_1}{l_2}M b{l'}{(\Delta, x : S \at l)}{(\boxit{T})} \\
    s_t = \ELIMTRM{l_1}{l_2}M b{(l \sqcup l')}{\Delta}{(\PI{l}{l'}xST)}{(\boxit{t})} \\
    s_s = \ELIMTRM{l_1}{l_2}M b{l}{\Delta}{S}{(\boxit{s})} \\
    \sigma = \Delta/g,S/U_S,T/U_T,t/u_t,s/u_s \\
    \delta = s_S/x_S,s_T/x_T,s_t/x_t,s_s/x_s}
  {\lttyequiv \metalevel {t_\tapp[l/\ell,l'/\ell',\sigma,\delta]}{\ELIMTRM{l_1}{l_2}M b{l'}\Delta{T_\tapp}{t'}}{M_\Trm[l'/\ell,\Delta/g,T_\tapp/U_T,t'/x_t]}{l_2}}
\end{mathpar}
Similar to above, we can do recursion on all sub-structures, including $S$ and $T$,
which are handled by the recursive principles for code of types. %
It is not only convenient to put $S$ and $T$ in the syntax of a function
application, but also necessary. %
If we look at $s_t$ and $s_s$, the recursive calls on the function and the argument,
we see that we must supply their types, i.e. $\PI{l}{l'}xST$ and $S$, respectively. %
This information, unfortunately, cannot be recovered, if we employed the more common
syntax of $t~s$. %
In practice, the $\Pi$ type can be filled in by a type inference algorithm when we do
not care, so it does not truly make the type theory more difficult, but rather enables
the recursion on code of function applications.

At this point, we conclude all rules for \delamlang. %
Next, we shall carefully define all syntactic operations and examine the syntactic
properties of \delamlang. %
Then we work out the semantics by following \Cref{sec:cv} and
\citet{abel_decidability_2017}, from which we conclude the convertibility problem of
\delamlang is decidable.

\subsection{A Note on Layer $\varlevel$ Rules}

To summarize, only the following rules can be indexed by layer $\varlevel$:
\begin{itemize}
\item the typing rule for local variables and its congruence;
\item the local substitution rules and their equivalence rules;
\item the typing rule for global variables and its congruence;
\item all conversion rules for terms and their equivalence.
\end{itemize}
In particular, we are not even obliged to include symmetry and transitivity, because
they can be derived from existing rules.
